\documentclass[10pt,a4paper]{amsart}
\usepackage[margin=19mm]{geometry}
\usepackage{amsmath,amssymb,mathtools}
\usepackage[scr=rsfs,cal=euler]{mathalfa}
\usepackage{xfrac}
\usepackage[unicode,hidelinks,bookmarks=false]{hyperref}
\newcommand{\ie}{i.e.\ }
\newcommand{\cf}{cf.\ }
\newcommand{\resp}{respectively\ }
\newcommand{\Subsection}{\subsection}
\providecommand{\nobreakdash}{\nobreak\hbox{-}}
\setlength{\emergencystretch}{2em}
\multlinegap0pt
\usepackage{amssymb}
\usepackage[shortlabels]{enumitem}
\usepackage{tikz}
\newtheorem{theorem}{Theorem}
\title{Quantization of Ricci Curvature in Information Geometry}
\author{Carlos C. Rodriguez}
\date{Source: arXiv:2603.10054v1 (CC BY 4.0)}
\begin{document}
\maketitle

\noindent\emph{Source and licence.} Quantization of Ricci Curvature in Information Geometry. Version arXiv:2603.10054v1 (CC BY 4.0), distributed by arXiv under the Creative Commons Attribution 4.0 International licence, http://creativecommons.org/licenses/by/4.0/, as recorded in the arXiv licence metadata for this exact version and shown on the abstract page https://arxiv.org/abs/2603.10054v1. Full licence notice: \texttt{licenses/arxiv-2603.10054-CC-BY-4.0.txt}.\par\smallskip

\noindent\textit{Editorial scope.} Counted unit: the article's theorem thm:gauss\_R (source0.tex lines 607--615). The article's complete original proof of it is delivered with it (source0.tex lines 617--643, one \texttt{proof} environment of 27 lines). No mathematics is written, repaired or re-derived: the statement and the proof are the article's own lines, in the article's own notation, and no retained line is substituted or re-flowed.  No further source-local context is needed: every cross-reference of every retained line resolves inside the two delivered spans.  Nothing was omitted from any retained span, so the package carries no omission record.  No retained line contains a citation, so the wrapper carries no bibliography and no bibliography entry.  No retained line contains a figure environment, an image-inclusion command or drawing code, so no image is delivered and the package declares no asset dependency.  Editorial formatting only, in addition: an AMS \texttt{amsart} preamble that restates the article's own declarations and macros used by the retained lines, namely the environments \texttt{theorem} on separate counters, together with the standard packages those lines need.  The retained environments print in this document as Theorem 1 in order of appearance, whereas the article prints the same items with the numbering of its own full document.  The complete native source of the article as distributed in the arXiv e-print for this exact version is bundled unchanged as \texttt{sources/196107/source0.tex}.
\begin{theorem}[Gaussian star curvatures]\label{thm:gauss_R}
For a Gaussian DAG with $n$ independent root nodes $X_1,\ldots,X_n$ and 
a single sink $X_{n+1}$ with all roots as parents 
(an $n$-root star, total dimension $d = 2n+1$), the Ricci scalar is 
the global constant
\begin{equation}\label{eq:gauss_R_formula}
  R = -\frac{n(n+3)}{2} = -\frac{(d+5)(d-1)}{8}.
\end{equation}
\end{theorem}

\begin{proof}
The parameter space has coordinates $(v_1,\ldots,v_n,a_1,\ldots,a_n,v_{n+1})$.  
After the log substitution $u_k = \tfrac{1}{\sqrt{2}}\ln v_k$, the metric becomes
$ds^2 = \sum_{k=1}^{n+1} du_k^2 + \sum_{i=1}^n e^{\sqrt{2}(u_i-u_{n+1})} da_i^2$.
This is a left-invariant metric on the solvable Lie group with orthonormal 
left-invariant frame $\{U_k,\, A_i\}$ (where $A_i = e^{-(u_i-u_{n+1})/\sqrt{2}}\partial_{a_i}$) 
and nonzero Lie brackets
\[
  [U_i, A_i] = -\tfrac{1}{\sqrt{2}} A_i, \qquad [U_{n+1}, A_i] = \tfrac{1}{\sqrt{2}} A_i.
\]
The Koszul formula gives $\nabla_{U_k}A_i = 0$ for all $k$ and 
$\nabla_{A_i}A_i = -\tfrac{1}{\sqrt{2}}U_i + \tfrac{1}{\sqrt{2}}U_{n+1}$.  
All sectional curvatures are computed from $K(X,Y) = \langle R(X,Y)Y, X\rangle$:
\begin{itemize}
\item $K(U_i, A_i) = -\tfrac{1}{2}$ for each $i=1,\ldots,n$ 
  (the $\nabla_{[U_i,A_i]}A_i$ term gives $-\tfrac{1}{2}U_i+\tfrac{1}{2}U_{n+1}$).
\item $K(U_{n+1}, A_i) = -\tfrac{1}{2}$ for each $i$.
\item $K(A_i, A_j) = -\tfrac{1}{2}$ for all $i \neq j$ 
  (siblings sharing the common sink; $U_{n+1}$ couples all $A_i$).
\item All other sectional curvatures are zero: 
  $K(U_i,U_j)=0$ (commuting), $K(U_i,A_j)=0$ for $i\neq j$ and $i\neq n+1$.
\end{itemize}
The total number of pairs with $K = -\tfrac{1}{2}$ is 
$n + n + \tbinom{n}{2} = 2n + \tfrac{n(n-1)}{2} = \tfrac{n(n+3)}{2}$.
Since $R = 2\sum_{i<j}K_{ij}$, we obtain $R = 2 \cdot \tfrac{n(n+3)}{2} \cdot (-\tfrac{1}{2}) = -\tfrac{n(n+3)}{2}$, 
which equals $-(d+5)(d-1)/8$ for $d=2n+1$.
\end{proof}

\end{document}
